\subsection{Open Mapping, Closed Graph, and Complemented Subspaces}

The Baire-category principle has three operator-theoretic consequences that should be read together: closed range gives uniform preimage estimates, open maps are detected at the unit ball, and closed graphs force boundedness.

\begin{theorem} \label{thm:closed-range-estimate} \hlt{Closed Range Estimate}\\
Let $\mathcal{X},\mathcal{Y}$ be Banach spaces, let $T\in\LL(\mathcal{X},\mathcal{Y})$, and set $\mathcal R\defeq T(\mathcal X)$. The following are equivalent.
\begin{enumerate}[label=(\roman*)]
\setlength{\itemsep}{0pt}
\item $\mathcal R$ is closed in $\mathcal Y$.
\item $\exists M>0$ such that $\forall y\in\mathcal R\ \exists x\in\mathcal X$ with $Tx=y$ and $\norm{x}\leq M\norm{y}$.
\end{enumerate}
\begin{proof}
\begin{enumerate}[label=(\roman*)]
\setlength{\itemsep}{0pt}
\item[$(ii)$]$\Rightarrow (i)$ Let $y\in\overline{\mathcal R}$. Choose $y_n\in\mathcal R$ with $y_n\to y$ and $\norm{y_n-y_{n-1}}\leq2^{-n}$ for $n\geq2$.\\
Pick $x_1,u_n\in\mathcal X$ with $Tx_1=y_1$, $Tu_n=y_n-y_{n-1}$, and $\norm{u_n}\leq M2^{-n}$.\\
Since $\sum\limits_{n\geq2}\norm{u_n}<\infty$, $x\defeq x_1+\sum\limits_{n\geq2}u_n$ exists and $Tx=\lim\limits_Ny_N=y$, so $y\in\mathcal R$.
\item[$(i)$]$\Rightarrow (ii)$ If $\mathcal R=\{\mathbf0\}$, take $M=1$. Otherwise $\mathcal R$ is Banach and $\mathcal R=\bigcup\limits_{n\geq1}n\,\overline{T(B_{\mathcal X})}^{\,\mathcal R}$.\\
By Baire and balanced convexity, $\exists c>0$ with $\overline{B}_{\mathcal R}\subseteq c\,\overline{T(B_{\mathcal X})}^{\,\mathcal R}$.\\
Thus each $0\neq r\in\mathcal R$ has $u\in\mathcal X$ with $\norm{u}\leq c\norm r$ and $\norm{r-Tu}\leq\frac{1}{2}\norm r$.\\
For $y\neq0$, iterate from $r_0=y$ to get $r_k=r_{k-1}-Tu_k$, $\norm{r_k}\leq2^{-k}\norm y$, and $\norm{u_k}\leq c2^{1-k}\norm y$.\\
Then $x\defeq\sum\limits_{k\geq1}u_k$ satisfies $Tx=y$ and $\norm{x}\leq2c\norm y$. For $y=0$, take $x=0$.
\end{enumerate}
\hlt{Structural Commentary}: Baire category turns surjectivity onto a closed range into uniform preimage control.\\
\hlt{Key Insight}: Closed range is equivalent to solving $Tx=y$ with a norm bound proportional to $\norm{y}$.
\end{proof}
\end{theorem}

\begin{proposition} \label{prop:relative-open-mapping-closed-range} \hlt{Relative Openness and Closed Range}\\
Let $\mathcal{X},\mathcal{Y}$ be Banach spaces, $T\in\LL(\mathcal{X},\mathcal{Y})$.\\
Then $T:\mathcal{X}\to T(\mathcal{X})$ is open, for subspace topology on range, iff $T(\mathcal{X})$ closed in $\mathcal{Y}$.
\begin{proof}
Let $\mathcal R\defeq T(\mathcal X)$.
\begin{enumerate}[label=\roman*.]
\setlength{\itemsep}{0pt}
\item[$(\Rightarrow)$] If $T:\mathcal X\to\mathcal R$ is open, then $\exists r>0$ with $B_{\mathcal R}(0,r)\subseteq T(B_{\mathcal X})$.\\
For $0\neq y\in\mathcal R$, choose $u\in B_{\mathcal X}$ with $Tu=\frac{r}{2\norm y}y$, and set $x\defeq\frac{2\norm y}{r}u$.\\
Then $Tx=y$ and $\norm{x}\leq\frac{2}{r}\norm y$; for $y=0$, take $x=0$. By Theorem \ref{thm:closed-range-estimate}, $\mathcal R$ is closed.
\item[$(\Leftarrow)$] If $\mathcal R$ is closed, then by Theorem \ref{thm:closed-range-estimate}, $\exists M>0$ such that every $y\in\mathcal R$ has $x\in\mathcal X$ with $Tx=y$ and $\norm{x}\leq M\norm y$. Let $U\subseteq\mathcal X$ be open and $x_0\in U$. Choose $\rho>0$ with $x_0+\rho B_{\mathcal X}\subseteq U$. If $z\in B_{\mathcal R}(Tx_0,\frac{\rho}{M})$, choose $h\in\mathcal X$ with $Th=z-Tx_0$ and $\norm h<\rho$; then $z=T(x_0+h)\in T(U)$. Thus $T(U)$ is open in $\mathcal R$.
\end{enumerate}
\hlt{Structural Commentary}: For linear maps, openness is detected at the unit ball around the origin.\\
\hlt{Key Insight}: Relative openness is the topological form of uniform preimage control.
\end{proof}
\end{proposition}

\begin{theorem} \label{thm:open-mapping-theorem} \hlt{Open Mapping Theorem}\\
Every surjective $T\in\LL(\mathcal{X},\mathcal{Y})$ between Banach spaces is open.
\begin{proof}
Since $T$ is surjective, $T(\mathcal X)=\mathcal Y$, $\mathcal Y$ is closed in itself. By Proposition \ref{prop:relative-open-mapping-closed-range}, $T:\mathcal X\to T(\mathcal X)=\mathcal Y$ open.\\
\hlt{Structural Commentary}: Surjectivity identifies the relative range with the codomain.\\
\hlt{Key Insight}: A Banach-space surjection has uniformly controlled preimages near the origin.
\end{proof}
\end{theorem}

\begin{corollary} \label{cor:banach-isomorphism-theorem} \hlt{Bounded Inverse Theorem}\\
If $T\in\LL(\mathcal{X},\mathcal{Y})$ is bijective between Banach spaces, then $T^{-1}\in\LL(\mathcal{Y},\mathcal{X})$.
\begin{proof}
Set $U\defeq T^{-1}:\mathcal Y\to\mathcal X$. By Theorem \ref{thm:open-mapping-theorem}, $T$ is open: if $O\subseteq\mathcal X$ is open, then $U^{-1}(O)=T(O)$ is open in $\mathcal Y$. Thus $U$ is continuous; since $U$ is linear, Theorem \ref{thm:bounded-linear-map-continuity-equivalences} gives $U\in\LL(\mathcal Y,\mathcal X)$.\\
\hlt{Structural Commentary}: Openness of the forward bijection is continuity of the inverse.\\
\hlt{Key Insight}: Banach completeness prevents a bounded linear bijection from having an unbounded inverse.
\end{proof}
\end{corollary}

\begin{corollary} \label{cor:banach-equivalent-norms-one-sided} \hlt{One-Sided Control of Banach Norms}\\
Let $\norm{\cdot}_1$ and $\norm{\cdot}_2$ be norms on $\mathcal X$ such that $(\mathcal X,\norm{\cdot}_1)$ and $(\mathcal X,\norm{\cdot}_2)$ are Banach.\\
Assume $\exists c>0$ such that $\norm{x}_2\leq c\norm{x}_1$ for all $x\in\mathcal X$.
\begin{enumerate}[label=(\roman*)]
\setlength{\itemsep}{0pt}
\item $\exists C>0$ such that $\norm{x}_1\leq C\norm{x}_2$ for all $x\in\mathcal X$.
\item $\norm{\cdot}_1$ and $\norm{\cdot}_2$ are equivalent.
\end{enumerate}
\begin{proof}
\begin{enumerate}[label=(\roman*)]
\setlength{\itemsep}{0pt}
\item Let $I:(\mathcal X,\norm{\cdot}_1)\to(\mathcal X,\norm{\cdot}_2)$ be the identity. Since $\norm{Ix}_2\leq c\norm{x}_1$, $I\in\LL$ and is bijective.\\
By Corollary \ref{cor:banach-isomorphism-theorem}, $I^{-1}\in\LL$, hence $\exists C>0$ such that $\norm{x}_1=\norm{I^{-1}x}_1\leq C\norm{x}_2$.
\item For all $x\in\mathcal X$, $\norm{x}_2\leq c\norm{x}_1$ by hypothesis and $\norm{x}_1\leq C\norm{x}_2$ by (i).\\
Hence the identity maps between the two normed structures are Lipschitz.
\end{enumerate}
\hlt{Structural Commentary}: Apply the bounded inverse theorem to the identity between the two normed structures.\\
\hlt{Key Insight}: Completeness upgrades one-sided norm control to two-sided equivalence.
\end{proof}
\end{corollary}

\begin{definition} \label{def:closed-operator-and-closed-complement} \hlt{Closed Operators and Closed Complements}\\
Let $\mathcal{X},\mathcal{Y}$ be normed spaces.
\begin{enumerate}[label=(\roman*)]
\setlength{\itemsep}{0pt}
\item \hlt{Closed operator}: $T:\mathcal{X}\to\mathcal{Y}$ is linear and $\forall(x_n)\subseteq\mathcal{X}$, $x_n\to x$ and $Tx_n\to y$ imply $Tx=y$.
\item \hlt{Closed linear complement}: $W\leq\mathcal{X}$ satisfies $\mathcal{X}=V\oplus W$ and is closed.
\end{enumerate}
\end{definition}

\begin{theorem} \label{thm:closed-graph-theorem} \hlt{Closed Graph Theorem}\\
Let $\mathcal{X},\mathcal{Y}$ be Banach spaces and let $T:\mathcal{X}\to\mathcal{Y}$ be linear. Set $\Gamma(T)\defeq\{(x,Tx):x\in\mathcal{X}\}$.
\begin{enumerate}[label=(\roman*)]
\setlength{\itemsep}{0pt}
\item $T$ is bounded iff $T$ is closed.
\item $T$ is closed iff $\Gamma(T)$ is closed in $\mathcal{X}\times\mathcal{Y}$.
\end{enumerate}
\begin{proof}
On $\mathcal X\times\mathcal Y$, use $\norm{(x,y)}\defeq\norm{x}+\norm{y}$.
\begin{enumerate}[label=(\roman*)]
\setlength{\itemsep}{0pt}
\item[$(i)$] ($\Rightarrow$) If $T$ is bounded, $x_n\to x$ and $Tx_n\to y$ imply $Tx_n\to Tx$; hence $y=Tx$, so $T$ is closed.
\item[$(i)$] ($\Leftarrow$) Assume $T$ is closed and set $\norm{x}_*\defeq\norm{x}+\norm{Tx}$. If $(x_n)$ is $\norm{\cdot}_*$-Cauchy, then $x_n\to x$ in $\mathcal X$ and $Tx_n\to y$ in $\mathcal Y$; by closedness, $y=Tx$, so $x_n\to x$ in $\norm{\cdot}_*$. Thus $(\mathcal X,\norm{\cdot}_*)$ is Banach.\\
Since $\norm{x}\leq\norm{x}_*$, by Corollary \ref{cor:banach-equivalent-norms-one-sided}, $\exists C>0$ such that $\norm{x}_*\leq C\norm{x}$; hence $\norm{Tx}\leq C\norm{x}$.
\item[$(ii)$] ($\Rightarrow$) If $T$ is closed and $(x_n,Tx_n)\to(x,y)$, then $x_n\to x$ and $Tx_n\to y$, so $y=Tx$ and $(x,y)\in\Gamma(T)$.
\item[$(ii)$] ($\Leftarrow$) If $\Gamma(T)$ is closed, $x_n\to x$, and $Tx_n\to y$, then $(x_n,Tx_n)\to(x,y)\in\Gamma(T)$; hence $y=Tx$.
\end{enumerate}
\hlt{Structural Commentary}: The graph norm turns closedness of the graph into completeness.\\
\hlt{Key Insight}: Graph-norm completeness and domination force equivalence with the original norm.
\end{proof}
\end{theorem}

\begin{theorem} \label{thm:closed-complements-continuous-projections} \hlt{Closed Complements and Continuous Projections}\\
Let $V$ be a closed subspace of a Banach space $\mathcal X$. The following are equivalent.
\begin{enumerate}[label=(\roman*)]
\setlength{\itemsep}{0pt}
\item $\exists W\leq\mathcal X$ closed such that $\mathcal X=V\oplus W$.
\item $\exists P\in\LL(\mathcal X)$ such that $P^2=P$ and $P(\mathcal X)=V$.
\end{enumerate}
\begin{proof}
\begin{enumerate}[label=\roman*.]
\setlength{\itemsep}{0pt}
\item[$(i)$]$\Rightarrow(ii)$ Let $\mathcal X=V\oplus W$ with $W$ closed, and define $P(v+w)\defeq v$. Then $P^2=P$ and $P(\mathcal X)=V$.\\
If $x_n\to x$ and $Px_n\to v$, then $v\in V$ and $(I-P)x_n\to x-v\in W$; by uniqueness in $V\oplus W$, $Px=v$.\\
Thus $P$ is closed, so $P\in\LL(\mathcal X)$ by Theorem \ref{thm:closed-graph-theorem}.
\item[$(ii)$]$\Rightarrow(i)$ Let $P\in\LL(\mathcal X)$ with $P^2=P$ and $P(\mathcal X)=V$. By Proposition \ref{prop:algebraic-projection-range-kernel-idempotent}, $\mathcal X=V\oplus\ker P$.\\
Since $\ker P=P^{-1}(\{\mathbf0\})$ and $P$ is continuous, $\ker P$ is closed.
\end{enumerate}
\hlt{Structural Commentary}: Closed complementary coordinates are exactly bounded projection coordinates.\\
\hlt{Key Insight}: The Closed Graph Theorem converts algebraic splitting into continuous splitting.
\end{proof}
\end{theorem}

\begin{theorem} \label{thm:complemented-range-is-closed} \hlt{Complemented Ranges Are Closed}\\
Let $T\in\LL(\mathcal{X},\mathcal{Y})$ act between Banach spaces. If $\operatorname{range}(T)$ has a closed linear complement in $\mathcal{Y}$, then it is closed.
\begin{proof} Set $\mathcal R\defeq T(\mathcal X)$, $\mathcal Y=\mathcal R\oplus Y_0$, $Y_0$ closed. Put $Z\defeq\mathcal X\times Y_0$, $\norm{(x,y_0)}_Z\defeq\norm{x}+\norm{y_0}$,\\
and $S:Z\to\mathcal Y$, $S(x,y_0)\defeq Tx+y_0$. Then $Z$ is Banach, $S\in\LL(Z,\mathcal Y)$, and $S(Z)=\mathcal Y$.\\
By Theorem \ref{thm:closed-range-estimate}, $\exists M>0$ such that $\forall y\in\mathcal Y\ \exists(x,y_0)\in Z$: $y=Tx+y_0$ and $\norm{x}+\norm{y_0}\leq M\norm{y}$.\\
If $y\in\mathcal R$, direct-sum uniqueness gives $y_0=0$; hence $\forall y\in\mathcal R\ \exists x\in\mathcal X$: $Tx=y$ and $\norm{x}\leq M\norm{y}$.\\
By Theorem \ref{thm:closed-range-estimate}, $\mathcal R$ is closed.\\
\hlt{Structural Commentary}: A closed complement transfers preimage control to the range summand.\\
\hlt{Key Insight}: Complementability supplies the missing uniform preimage estimate for closed range.
\end{proof}
\end{theorem}

\begin{corollary} \label{cor:finite-codimensional-operator-range-closed} \hlt{Finite-Codimensional Operator Ranges Are Closed}\\
Let $\mathcal{X},\mathcal{Y}$ be Banach spaces and let $T\in\LL(\mathcal{X},\mathcal{Y})$. If $\operatorname{codim}_{\mathcal{Y}}T(\mathcal{X})<\infty$, then $T(\mathcal{X})$ is closed.
\begin{proof}
Set $\mathcal R\defeq T(\mathcal X)$. Since $\operatorname{codim}_{\mathcal Y}\mathcal R<\infty$, choose $Y_0\leq\mathcal Y$ with $\mathcal Y=\mathcal R\oplus Y_0$ and $\dim Y_0<\infty$.\\
By Theorem \ref{thm:finite-dimensional-normed-spaces}, $Y_0$ is closed; by Theorem \ref{thm:complemented-range-is-closed}, $\mathcal R$ is closed.\\
\hlt{Structural Commentary}: Finite codimension produces a finite-dimensional closed complement.\\
\hlt{Key Insight}: Only finitely many missing codomain directions force the operator range to be closed.
\end{proof}
\end{corollary}
